\section{Related Work}
\label{sec:related}

% 迁移自 docs/related_work.md；bib 条目已全部核实（2026-08）。

\subsection{ML-Driven Traffic Engineering for WANs}

Traffic engineering has long been formulated as a constrained optimization
problem~\cite{fortz2000,wang2006} solved by linear programming (LP), whose
runtime scales poorly with network size. Kumar \emph{et al.} proposed
SMORE~\cite{smore}, which decouples robust oblivious path selection
from online rate adaptation, while Abuzaid \emph{et al.}~\cite{ncflow} and
Narayanan \emph{et al.}~\cite{pop} accelerate LP by problem decomposition,
trading allocation quality for speed. A recent line of work replaces online
solvers with learned models~\cite{valadarsky2017}: Perry \emph{et al.}
proposed DOTE~\cite{dote}, which directly maps historical traffic matrices to
split ratios with stochastic optimization; Xu \emph{et al.} proposed
Teal~\cite{teal}, combining a flow-centric graph neural network (GNN),
multi-agent reinforcement learning (RL), and alternating direction method of
multipliers (ADMM) fine-tuning to achieve near-optimal allocations with
orders-of-magnitude speedups; and Liu \emph{et al.} proposed
FIGRET~\cite{figret}, which enhances robustness against traffic uncertainty at
per-flow granularity. To survive topology changes without retraining, AlQiam
\emph{et al.} proposed HARP~\cite{harp}, enforcing invariances to node
permutation and tunnel reordering so that it transfers across topologies unseen
in training; Fan \emph{et al.} proposed Aether~\cite{aether2025}, pursuing the
same generality with elastic multi-agent graph transformers, and Ye
\emph{et al.}~\cite{ye2025} trained path-based GNNs for robust distributed TE
under link failures. More recently, Wu \emph{et al.}~\cite{netllm} explored
large language models as general-purpose traffic planners. However, all of
these systems assume the \emph{complete} traffic matrix is available at
decision time, an assumption that rarely holds in large-scale satellite
networks, where network-wide measurement is prohibitively expensive.

\subsection{TE and Routing for Satellite Networks}

LEO mega-constellations pose unique challenges for TE: topologies change
every few tens of milliseconds as satellites move, and control decisions
must be produced within stringent latency
budgets~\cite{handley2018,bhattacherjee2019,kassing2020}. Classical
approaches
absorb mobility through topology abstractions; the virtual node model of Lu
\emph{et al.}~\cite{vn2013} binds logical nodes to earth-fixed footprints so
that upper-layer routing observes a quasi-static grid. Optimization-based
satellite TE has explored segment routing~\cite{hu2022}, distributed
clustering that optimizes elephant flows separately~\cite{wei2024}, and
supervised learning for distributed split-ratio prediction~\cite{flexsate2024}.
Wu \emph{et al.} proposed SaTE~\cite{sate}, the
state-of-the-art learning-based TE for satellite constellations: it
formulates a heterogeneous graph over satellites, demands, paths, and links,
prunes zero-traffic relations to make training tractable on a single GPU,
and selects representative topologies via Graph2Vec and determinantal point
process (DPP) sampling so that
one trained model generalizes across topology snapshots with millisecond
inference. In parallel, Yang \emph{et al.} proposed FNC~\cite{fnc,fnc-infocom},
which targets dynamic traffic
allocation for earth-observation constellations, replacing RL with
end-to-end imitation learning and substituting iterative constraint solvers
with normalization-based feasibility layers. Notably, FNC's imitation
learning requires expert allocation labels computed from complete training
matrices. Our offline training also uses complete traffic to evaluate rewards,
but does not require expert allocation labels; the distinction is not the
availability of training-time ground truth. Both SaTE and FNC, like their WAN counterparts,
assume the traffic matrix is fully observable at decision time. Unlike SaTE
and FNC, \sysname{} performs TE from sparse observations, down to a $2\%$
observation ratio (\cref{tab:related}).

\subsection{TE with Incomplete Traffic Observations}

Network-wide traffic measurement is expensive: collecting a complete
traffic matrix requires instrumenting all ingress nodes at fine time
granularity, which is particularly onerous on resource-constrained
satellite platforms. A classical workaround is a two-stage pipeline:
first estimate or complete the traffic matrix via network tomography
(Vardi~\cite{vardi1996}; Zhang \emph{et al.}~\cite{zhang2003}), matrix
completion, or compressive sensing~\cite{zhang2009}, then optimize routing on
the estimate. However, reconstruction errors propagate to and degrade the
downstream allocation, since a single reconstructed matrix is handed to a
separate optimizer. Guo \emph{et al.} proposed TEST~\cite{test}, which, to
our knowledge, is the first \emph{end-to-end} approach to learn routing
policies directly from sparsely measured traffic matrices rather than first
reconstructing them; it integrates a Transformer over historical sparse
sequences with a graph convolutional network (GCN) over the topology, fills
unobserved entries with zeros, compensates with synthesized training
augmentation, and targets terrestrial networks with static topologies. In the
graph learning literature, Huo \emph{et al.}~\cite{t2gnn} addressed message
passing under missing node features via teacher-student distillation, but such
techniques have not been applied to network TE. No prior work, to our
knowledge, performs TE from sparse traffic observations in satellite networks,
where the difficulty is compounded by time-varying topologies and drifting
demand sets. Unlike TEST, which fills unobserved entries with zeros and
targets static terrestrial networks, \sysname{} replaces missing entries with
learnable mask embeddings, gates their message passing, and handles LEO
dynamics---time-varying topologies and drifting demand-pair sets---via demand
pruning with size-invariant weights, enabling train-once, zero-retraining
deployment; it further uses a source-neighbor estimate for optional
congestion-aware refinement without consulting the unobserved ground-truth
matrix (\cref{tab:related}).

\begin{table}[t]
    \caption{Comparison with the related works.}
    \label{tab:related}
    \centering
    \footnotesize
    \setlength{\tabcolsep}{3pt}
    \begin{tabular}{@{}lccccc@{}}
        \toprule
        & Teal~\cite{teal} & SaTE~\cite{sate} & TEST~\cite{test}
        & FNC~\cite{fnc} & \textbf{Ours} \\
        \midrule
        Satellite scenario      & \xmark & \cmark & \xmark & \cmark & \cmark \\
        Dynamic topology        & \xmark & \cmark & \xmark & \cmark & \cmark \\
        Sparse observation      & \xmark & \xmark & \cmark & \xmark & \cmark \\
        Temporal modeling       & \xmark & \xmark & \cmark & \xmark & \cmark \\
        Mask-aware model        & \xmark & \xmark & \xmark & \xmark & \cmark \\
        Neighbor estimate        & \xmark & \xmark & \xmark & \xmark & \cmark \\
        Demand-set drift        & \xmark & partial & \xmark & \xmark & \cmark \\
        \bottomrule
    \end{tabular}
\end{table}
